\chapter{Why Language Keeps Changing}
\section{A Family Letter You Cannot Quite Read}
First pick up a letter.

Not email, but real paper: yellowed, brittle at the corners, folds so deep they nearly split. You found it while sorting your grandmother's belongings after her death. In 1983 she wrote to your mother, then at university. The handwriting is neat; the opening says ``Seeing these words is like seeing me.'' One passage reads:

\begin{quote}
We still have enough grain ration coupons. Last month your father managed through someone to get an industrial-goods coupon. He wants to save it to buy you a Dacron shirt, cool to wear when the weather gets hot. A young person has just been assigned to our work unit, very ambitious, a Youth League member\ldots{}
\end{quote}
You read it three times. You recognize every character, genuinely every one, yet the passage seems behind frosted glass. Grain ration coupons you have seen in history books: vouchers for buying food. But what was an industrial-goods coupon? Was díquèliáng, Dacron, a fabric, and why does its Chinese name sound like the start of a compliment, ``indeed excellent''? Why begin praising someone by calling them ``ambitious'' and ``a Youth League member''?

Stranger still is the tone. The grandmother you remember loved jokes, yet this letter is solemn, without one. Your mother explains: that was how people wrote then. Letters might be read aloud publicly, so everyone wrote as if delivering a speech.

Suddenly you realize: only forty years separate you from this letter, two generations of mothers and daughters. The language is the same ``Chinese,'' the same square characters, yet you already need a translator. Not because you cannot read the characters, but because their world has vanished. Grain coupons disappeared in 1993; Dacron faded as department stores became ``shopping centers''; even the letter's tone vanished into WeChat stickers.

Words have lifespans. The letter is amber preserving 1983 Chinese. Outside the amber, Chinese has long since flowed elsewhere.

This brings us to our question: why does language change, and how fast? Is change ``decline,'' as adults often claim when saying young people's speech gets worse and worse, or something else? The letter is our starting point, but answering requires a more distant scene.

\section{A Graffiti Wall Beneath Volcanic Ash}
Come with me to Pompeii in 79 CE.

At Mount Vesuvius's foot in southern Italy, Pompeii was a lively Roman port city with ten or twenty thousand residents, taverns, baths, an arena, and bakeries. That August, the volcano erupted and buried it alive beneath meters of ash. More than seventeen centuries later, archaeologists excavated it. The ash had sealed everything like cement, including the writing on walls.

Stand on a Pompeian street. It is afternoon, the sun fierce; wheels have worn two deep grooves in the paving stones. A tavern stands beside you, its large counter jars warming wine. Bread and fish sauce scent the air. Its plaster wall is dense with writing, not official inscriptions but ordinary people's casual charcoal marks and scratches made with nails.

What is on it? Campaign slogans: roughly, ``Support so-and-so for municipal office.'' Love notes: one man writes his beloved's name; another crosses it out and supplies his own. Accounts: ``On this date, so-and-so owes me wine money.'' One line later readers remember amounts to ``Innkeeper, I demand repayment for your watered wine.'' Others complain, swear, and draw cartoons, like walls in every era.

These scribbles are a linguist's gold mine because they record \textbf{ordinary people's spoken Latin}, not Cicero's speeches. Roman schools taught standard written Latin with rigorous grammar and complex endings: a noun changed its tail six ways according to its sentence role. Pompeian graffiti gives everyday speech away. Writers repeatedly make ``mistakes,'' often dropping final m and confusing vowels. Spelling mistakes leak pronunciation: people often omit a letter because their mouths have already stopped saying it.

Thus in 79 CE, ``correct Latin'' and ``street Latin'' already wore different faces. Schools taught six noun endings; everyday speakers had simplified several. Vowels carefully distinguished in writing had merged in speech. Centuries later, the Roman Empire collapsed, roads deteriorated, and local spoken Latins drifted onward in separate villages and ports. Spain's port Latin, northern France's farm Latin, Italy's mountain-town Latin grew increasingly unintelligible to one another. A millennium later they had new names: Spanish, French, Italian, Portuguese, Romanian.

Today a Spaniard says agua, ``water,'' and a French speaker eau, pronounced like ``oh.'' They look unrelated, yet both trace back to Latin aqua. One river flows until its branches no longer recognize one another.

Remember this graffiti wall. It tells us that linguistic change begins not in ``getting worse'' but in \textbf{use}. Every ordinary person's casual mark casts a tiny vote for change.

\section{Where Does Language Change?}
Before judging change good or bad, identify where it occurs. Linguists separate several levels. Each changes, but at different speeds.

\textbf{First: pronunciation changes.}

Chinese offers a famous example. Earlier Chinese had an ``entering tone'': short syllables ending abruptly in a stop consonant. In today's Mandarin, one, seven, eight, and ten have no such difference in duration from other syllables. Ask a Cantonese speaker from Guangzhou: jat, cat, baat, sap each ends sharply. Cantonese retained the entering tone; Mandarin lost it centuries ago. Nobody ``decided'' to abolish it. Around Beijing, mouths wore it smooth generation after generation.

Pompeii's lost final m is the same kind of event. Pronunciation is a shoe sole: daily walking wears it down.

\textbf{Second: meanings change.}

Classical Chinese zǒu meant ``run.'' In zǒumǎ guānhuā, ``view flowers from a running horse,'' the horse runs, not strolls. The word for walking was xíng. Tāng meant scalding water, making ``enter boiling water and tread through fire'' a fearsome idiom. If it meant today's seaweed-and-egg soup, the phrase would lose its terror. Meanings drift like boats: some broaden, some narrow, some change harbors entirely. Open any ``ordinary word,'' and you may find a bizarre résumé.

\textbf{Third: grammar changes.}

The Analects has wèi zhī yǒu yě, ``there has been no such thing.'' Notice the order: in a negative clause, pronoun zhī precedes the verb. That was the pre-Qin rule; today anyone saying it sounds theatrical. Earlier Chinese did not require classifiers as modern yī běn shū, ``one classifier book,'' and sān pǐ mǎ, ``three classifier horses,'' do; ``horses three pǐ'' was normal. The whole array of particles zhī, hū, zhě, yě has likewise retired from everyday use. Grammar looks most stable, but moves too, slowly as a glacier.

Put these layers together, and our real conflict appears:

\textbf{Almost every era has people announcing that language is declining; yet each era's standard is precisely the product of the previous era's ``decline.''}

Today's standard French is ``distorted Latin.'' Today's Mandarin is Middle Chinese with the entering tone worn away. Every rule in your ``standard written language'' rests on some former act of ``recklessness.'' If change equals decline, two thousand years should have reduced human language to monkey calls. Yet you can use it to read physics problems, write love letters, and program computers.

Why, then, is the intuition of decline so stubborn? How likely is it to be wrong? First hear what each side says.

\section{Purists and the Living-Language Camp}
Two camps have argued about change for centuries. Let us hear both before taking sides.

\textbf{Purists: language needs guardians.}

The French are especially serious. Founded in 1635, the Académie française guards French ``purity,'' compiles an authoritative dictionary, and judges which words qualify as French. English loans are treated almost as enemies. Do not simply use computer; use their coined ordinateur. Prefer courriel to e-mail. China has echoes: recurring campaigns to ``protect Chinese purity,'' official requests that broadcasters limit foreign abbreviations like NBA and GDP, teachers distressed by juéjuézi or yyds in essays.

Before laughing, state their full reasons. We practice this repeatedly: strengthen the other side's argument before deciding whether to refute it. Purists have at least four substantial reasons.

First, \textbf{clarity is common property}. Language is everyone's road. If meanings change at every use, signposts become unreliable. In laws, contracts, and medicine instructions, ambiguity can cost lives.

Second, \textbf{generations need bridges}. Change too quickly and grandparents cannot understand grandchildren; society shares fewer topics. Needing help to understand your grandmother's letter is no pleasant feeling.

Third, \textbf{some new words really are worse}. Replace a precise old word with a vague new one and you incur a net loss. If internet usage flattens several distinctions into one fuzzy compliment, expressive capacity genuinely shrinks.

Fourth, \textbf{language holds memory}. Lose tāng's older meaning, hot water, and ``enter boiling water and tread through fire'' loses force. Old words and meanings preserve millennia of collective life.

Together these make a respectable position: language may change, but when change is too rapid or casual, someone has a right to apply the brakes.

\textbf{The living-language camp: language is for use, not worship.}

The other side has strong reasons too. Language's first function is communication now, not museum display. Whether a word is ``good'' has one judge: its countless users' mouths. Nearly four centuries of the French Academy have not stopped French people saying le weekend. Official documents cannot keep ``pay the bill'' and ``take a taxi'' expressions out of Chinese dictionaries. No institution has ever successfully frozen a living language. Only dead languages, like Latin, freeze: no longer used to shop or quarrel, they become ``pure,'' and their purity is bound up with death.

This camp also holds our earlier counterexample as a decisive weapon: \textbf{the ``orthodox French'' purists defend is yesterday's ``distorted Latin.''} Revive a Roman teacher from 79 CE and let them hear French: they would despair at lost endings, altered sounds, total impropriety. French does not consider itself degenerate Latin, but proper French. Likewise, today's critics of internet language speak ``standard Chinese'' that would sound distorted to Qing-era ears. Each generation's purists stand on the ruins of the previous generation's ``decline'' and declare the ruins a temple.

This does not make purists wholly wrong. ``Language needs brakes'' and ``language inevitably flows'' can be two gears of one system. The real dispute is who holds the brake: academies, dictionaries, and official documents, or everyone currently speaking? History's vote is clear: speakers do. But that answer is too coarse. We need a tool to inspect change's finer structure.

\section{Thinking Bridge: Language Paternity Testing}
The tool is the \textbf{historical-comparative method}, honed by nineteenth-century linguists. Like a biologist's DNA test, it establishes languages' relationships and reconstructs long-vanished ancestors.

The story often begins in 1786. British judge William Jones, working in Calcutta, studied India's classical Sanskrit in his spare time. He became increasingly astonished: its similarities to ancient Greek and Latin were too extensive for chance. Not a few coincidentally similar words, but a systematic pattern. He publicly proposed that all three probably descended from a vanished common ancestor. Later work confirmed his intuition and greatly expanded the family: dozens of languages from India through Iran to Europe, including English, German, Russian, Persian, and Hindi, belong to \textbf{Indo-European}.

What makes similarity ``systematic''? Compare Latin pater and English father; Latin piscis and English fish; Latin ped-, the foot root behind English pedal, and English foot. The pattern is clear: Latin p corresponds consistently to English f. Pater/father, piscis/fish, ped-/foot align pair after pair. A few chance borrowings would not produce such regularity. A correspondence orderly enough to be a law, Grimm's law, has only one explanation: \textbf{they were once one language; after separation, each drifted by its own rules, leaving traces like growth rings in every word.}

The method's central operations are:

\begin{itemize}
\item Gather basic vocabulary least likely to be borrowed: numbers, body parts, kinship terms, water, fire, sun, moon.
\item Compare pronunciations across languages, looking for \textbf{series of regular correspondences}, not isolated resemblances.
\item Work backward from those correspondences: what was the ancestral sound probably like, and what form did the word have?
\end{itemize}
Using this method, linguists reconstructed Proto-Indo-European, a language with no written record that nobody has directly encountered, inferred solely from descendants' growth rings. The rough picture is that a community spoke it about six thousand years ago; descendants carried it toward Europe and South Asia, separating farther until one language became dozens.

Chinese has a family too. Basic-vocabulary comparisons lead linguists to relate it to Tibetan, Burmese, and others within \textbf{Sino-Tibetan}. This field is much younger than Indo-European research, and many details remain disputed. That itself is a reminder: relationship tests differ in strength. Strong evidence supports established conclusions; weaker evidence must carry a question mark. The method's discipline is: \textbf{without law-like correspondences, better withhold the resemblance claim.}

This tool gives you new eyes. Whether languages resemble one another is decided by correspondence tables, not feelings. ``They were once a family'' becomes not a myth, as many peoples' common-origin legends are, but a checkable problem.

\section{Two Excavated Samples of ``Standard Language''}
Now read two short primary texts, Chinese and European, each its culture's ``standard language'' two millennia ago.

First China. In ``On Government'' in the Analects, Confucius says:

\begin{quote}
Know what you know; acknowledge what you do not know. This is wisdom.
\end{quote}
Still quoted today, the sentence's Chinese grammar and vocabulary bear time's marks. Word by word: knowing is knowing, not knowing is not knowing; shì here is not today's linking verb ``is,'' but the demonstrative ``this.'' The final character normally read zhī, ``know,'' is read zhì, ``wisdom.'' The whole means that honestly distinguishing knowledge from ignorance is a wise attitude toward knowledge.

To reach a modern reader, the brief sentence needs two translations: shì changes meaning (this $\rightarrow$ is), and zhī changes reading and function (know $\rightarrow$ wisdom). This is live evidence of semantic and phonetic drift across more than two thousand years. The Analects was conversational recorded speech in its own era. Today it is annotated ``classical Chinese,'' not because it grew more profound, but because language moved on while it stayed put.

Now Europe. Near the beginning of Matthew chapter six in the Latin Vulgate Bible comes the famous prayer's opening:

\begin{quote}
Pater noster, qui es in caelis, sanctificetur nomen tuum.
\end{quote}
In essence: Our Father in heaven, may your name be held holy. Take this Latin to today's speakers and hear echoes: Italian Padre nostro, che sei nei cieli; French Notre Père, qui es aux cieux; Spanish Padre nuestro, que estás en los cielos. Pater $\rightarrow$ padre $\rightarrow$ père; noster $\rightarrow$ nostro $\rightarrow$ nuestro $\rightarrow$ notre. The same ``standard language,'' worn by three sets of mouths in three places, takes three forms. Each claims authenticity, and each is authentic.

Together the passages show that \textbf{a standard language is neither a beginning nor an end, but a photograph of the river in a particular year}. Confucius's Chinese was not the ``original authentic product''; it descended from earlier Chinese. Latin was not Romance's ``paradise before decline''; it was itself a branch flowing from Proto-Indo-European. The photograph may be sharp, but the river keeps moving. Declare one photograph the only correct language, and you protect not language's essence but a year.

Again distinguish four things: what happened (documented changes in meaning and sound); why (next section); how contemporaries understood it (Confucius did not know he spoke ``ancient Chinese,'' Romans did not know they spoke ``future French''---nobody lives in a ``transition,'' every era thinks itself final); and how we judge it (whether change is decline is evaluative, not factual). Mixing these is the topic's commonest mistake.

\section{One River, Three Accounts of Its Flow}
The facts are agreed: language keeps changing. Explanations compete over \textbf{why}. Here are three seriously argued accounts, each holding some truth and leaving something unreachable.

\textbf{Explanation one: economy of effort---change comes from laziness.}

The simplest account: people prefer easier speech. Speech organs, like all tools, take the convenient route. Latin aqua has two fully pronounced syllables; French eau reduces them to one sound, wearing away consonants and merging vowels. Nobody pronounces every letter of English Wednesday. Beijing speech often weakens the second syllable of dòufu, ``tofu,'' until its vowel nearly vanishes. Language is a mountain path worn by millions: it shifts toward easier walking.

This account is powerful and vivid, but has an obvious gap: \textbf{it cannot explain change's timing and location}. If everyone seeks ease, why did northern speech lose entering tones while Cantonese preserved them? Why does a sound change happen this century rather than last? Pure economy should make every language slide toward the same ``easiest language,'' yet they diverge. Laziness is universal, changes local; universality alone cannot explain locality.

\textbf{Explanation two: identity---change is a young person's ear piercing.}

Twentieth-century sociolinguists changed viewpoints. They entered specific communities and counted who led pronunciation changes. In a famous island study, American linguist William Labov found that islanders, especially young people intending to stay and resist tourists' culture, strengthened and exaggerated particular older local pronunciations. Change may begin not with laziness but with a \textbf{declaration}: I differ from you, I belong here, I am young.

This illuminates many things. Why does every adolescent generation have slang that expires once parents learn it? Because it marks identity, whose point is that parents remain outside. Why do online buzzwords live ever shorter lives? Because the identity of ``someone still using that word'' updates so quickly. Change gains a social engine: a word or pronunciation often spreads not because it works better but because it marks ``us.''

The limitation: identity explains bursts of new words and accent choices better than the \textbf{centuries-long, law-like regularity} found by comparative linguistics. Grimm's p-to-f change affects entire word sets and populations with centuries of consistency. Nobody sustains an ``ear piercing'' for centuries. Social motives ignite a spark, not tectonic movement.

\textbf{Explanation three: rupture---history drives change.}

This account zooms out furthest. Major linguistic shifts often coincide with historical breaks. Roman roads and administration tied local Latins together; imperial collapse cut the ropes and freed dialects to separate. Romance's great divergence occurred in the following centuries. English is more dramatic. In 1066, Norman French conquerors took England; French became the language of court and law, while English was relegated to kitchens and fields. When English reemerged three centuries later, it had a different body: greatly simplified grammar and a flood of French words. Field animals kept English names, cow and pig; dishes on aristocratic tables took French-derived beef and pork. The mark remains in dictionaries. Chinese endured successive impacts too: Buddhism brought terms such as ``world'' and ``instant''; modern Western learning brought ``science'' and ``democracy,'' many passing through Japanese translation before returning.

Rupture explains large vocabulary renewals and grammatical turns. Its weakness: shocks explain ``why then,'' but not the steady internal drift that continues without external events. Language moves in peaceful times too.

The three accounts are three hydrologies: economy studies riverbed material, identity the eddies, rupture the floods that redirect the channel. All operate in a real river. Together they reject one claim: \textbf{none equates change with decline}. Economy is engineering optimization, not regression; identity marking is social instinct; rebuilding after rupture is living things answering a living world.

\section{Further Challenge: Grammar from Fragments}
Now enter the chapter's weightiest theoretical scene. If change is not decline, can its apparent opposite exist: \textbf{creating a language from scratch}?

First two concepts. A \textbf{pidgin} arises when adults without a common language must interact daily: laborers abducted or traded from different continents onto plantations, sailors from different countries in ports. They assemble an emergency language, vocabulary pieced together, grammar minimal, without tense, endings, or subordinate clauses. ``I went to the market to buy fish yesterday'' becomes something like ``I go market buy fish yesterday.'' A pidgin is nobody's mother tongue: an adult tool, rough and makeshift, enough to get by.

A \textbf{creole} emerges in the next generation. Children born into a pidgin environment learn this ``fragmented language'' natively, and something extraordinary happens. Linguists report that within a single generation, creoles develop the equipment pidgins lack: systematic tense markers, nested clauses, consistent pronunciation. Haitian Creole, largely French in vocabulary, is now the native language of over ten million, with literature, news, and a constitution. Papua New Guinea's Tok Pisin began as a plantation pidgin and is now an official language. Fragments go in; a complete language comes out.

How to explain it? American linguist Derek Bickerton proposed the bold \textbf{language bioprogram hypothesis}: children's brains contain a grammatical blueprint. However fragmented the input, they fill missing structures according to it; creoles reveal the blueprint on an unfinished building. The hypothesis provoked long debate. Later research exposed weaknesses: creole grammars have many specific links to their African and European source languages, so need not all come from one universal template. ``Built in one generation'' may also compress decades of gradual change into a legend. Treat it as a \textbf{powerful, unfinished explanation}, not a verdict.

Even discounting Bickerton, the central fact remains: \textbf{children's linguistic creativity far exceeds what decline theories imagine}. If change meant generational wear, children in pidgin environments should learn less than their parents. The reverse happens. They are the one group repeatedly observed ``creating languages'' in human history. Change is construction as well as erosion.

The scene has a second half in the opposite direction: \textbf{language loss}. Of roughly seven thousand languages worldwide today, more than half have few speakers and are no longer learned by children. Estimates from UNESCO and similar institutions are often cited: current trends may carry away a substantial number, some say nearly half, with their last speakers this century. Linguists say that, at regular intervals, another language falls forever silent somewhere in the world.

Notice that most language loss is not ``natural death.'' Indigenous Australian children were sent in groups to boarding schools and punished for speaking their languages. Similar things occurred in the Americas, northern Europe, and Siberia. Many languages were abandoned under force. Saying a language ``disappeared'' sounds like a leaf falling; often people were first forced into silence. Some communities now work in reverse. Māori communities in New Zealand founded ``language nest'' preschools, immersing young children in Māori. Hebrew offers a still more extreme story: silent in ordinary speech for over a millennium, surviving in religious texts, it was brought back to the street by a generation from the late nineteenth century onward. Millions now shop and quarrel in it. If a language ``dead'' for a thousand years can revive, loss is no physical law but the accumulated result of choices. That also makes every choice someone's responsibility.

By now, ``language is declining'' no longer stands. But a harder question emerges: if change is normal and construction instinctive, how should we view today's words, memes, and abbreviations?

\section{Language's Future in Your Input Method}
Return to today. Every old question in this chapter is streaming live in your pocket.

\textbf{New words, loans, and abbreviations enter and leave Chinese faster than ever.} Your input method contains yyds, initials for ``eternally the greatest''; xswl, ``laughing to death''; emo; involution; and ``defenses breached.'' The internet did not invent abbreviation. Běijīng Dàxué becomes Běidà, ``Peking University''; ``Olympic Games'' becomes a shorter Chinese form without complaint. English radar began as radio detection and ranging; laser likewise is an acronym. Abbreviation is an old compression technology. The internet merely accelerates compression and obsolescence to the limit: a word can be born and passé within a summer.

Borrowing is similar. You may call it ``worshiping foreign things,'' yet your own mouth is full of unnoticed loans. Pútao, ``grape,'' came from the Western Regions in the Han era; shī in ``lion'' traces to Old Iranian. Sofa, coffee, and chocolate arrived in modern times; logic and humor more recently. Economy, science, philosophy, and many others circulated through Japanese translation. Chinese has borrowed for over two thousand years. Long use gives newcomers residency, and everyone forgets their foreign origins. Opponents of borrowing use dictionaries packed with loans.

\textbf{You already live amid official pronunciation review.} Many do not know that older dictionaries pronounced dāibǎn, ``rigid,'' as áibǎn. A national review of variant pronunciations standardized dāibǎn in 1985. Teachers who had corrected students for half a lifetime suddenly found their favored version was the ``old wrong one.'' This was not somebody's whim, but official recognition of change. The majority's mouths moved first; the standard followed with a stamp. River first, map afterward, as always.

\textbf{Internet language is the newest Pompeian graffiti.} Linguists study it just as they study Pompeii: examine ordinary people's casual writing, not textbooks. Your group-chat ``hahahaha'' and ``haha'' are not the same laugh. Punctuation, character count, and emojis become new tone particles; Chinese grows new ``grammatical organs'' on screens. Imagine linguists two millennia from now excavating today's chat servers. They would be as delighted as we are by Pompeian walls: ``So this was colloquial Chinese in the early twenty-first century!'' They would not conclude ``Chinese had declined,'' but ``Look, change was happening, and we caught it.''

\textbf{A wholly new variable has entered: machines speak too.} AI models train on billions of people's texts, add their writing to the internet, and are then read by the next generation. Language has acquired a ``human-machine feedback'' loop for the first time. Nobody knows where it will take Chinese. That is the unfolding scene this chapter leaves you.

\section{Over to You: Should Language Need an ID Card?}
Return to the family letter, Pompeii's wall, and your current input method.

You now have all the material. Sounds wear, meanings drift, grammar moves. Comparative methods establish that separated languages were siblings. Economy, identity, and rupture drive together. Fragmented language can become complete in children's mouths; complete language can die in a generation's silence. Today's defended ``purity'' was yesterday's denounced ``decline.''

So comes the final question:

\textbf{Suppose your city proposes a ``language protection committee,'' authorized to decide which new words count as standard Chinese and enter textbooks, which are pollution to remove from media and exams, whether broadcasts may use foreign abbreviations, and whether student essays lose marks for internet buzzwords. Vote: support the committee or oppose it?}

Give your answer and reasons. Before concluding, weigh both sides again.

Support has real grounds: language is common property; ambiguity and generational gaps have costs; some new words are less precise; language holds millennia of memory. And you know that once a usage disappears, nobody remains to vote for it.

Opposition has real grounds too: no committee has frozen a living language; frozen ones are dead. The standards committees defend grew wild before such oversight, and ``errors'' are often a future not yet officially recognized.

Deeper lies a harder question: who sits on the committee? Does its ``standard'' happen to match its members' class, age, and region? Whoever issues language an ID card really holds a stamp classifying \textbf{people}. Those whose speech is ``standard'' count as respectable. Pompeian tavern scribblers, Gallic farmers speaking ``degraded Latin,'' and your relative with ``nonstandard Mandarin'' all know the experience.

There is no standard answer. One thing is certain: whichever vote you cast, your manner of voting---word choice, sentence form, whom you consult---casts a real vote for Chinese's future. Committees have never decided that future. You do, along with everyone speaking now. When linguists excavate our era two millennia hence, you too will be on that wall.

May the mark you leave be one you thought through.
